\subsection{Theme Synthesis — 競争軸の分化}

4月の断面から見えるのは、frontier modelの競争が消えたのではなく、その意味が広がったことだ。model capabilityは依然として中心にあるが、実際の利用条件はAPI lifecycle、互換interface、multimodal surface、agent接続、open-source libraryの追随によって形作られる。5月以降に実行環境全体の競争が前景化する前に、4月にはすでにその構成要素が別々の方向へ伸び始めていた。本号のfrontier比較は、その分岐点を捉えるためのものであり、vendor benchmarkを単一ランキングへ変換するためのものではない。\autocite{src-6ddeedd7665f,src-9156c37fb1c2,src-a24fd97eb3d0,src-dfe7e0c6ff39,src-e547d0d8cbe3,src-eac8dfbeb78c}


\begin{claimboundary}
Claim Boundary: first-partyのrelease pageや公式feedは、公開chronologyとsource ownerが述べたscopeを確認する一次資料である。一方、「より賢い」「高性能」といったcapability・benchmark・safety・performanceの評価は、それだけで独立再現済みの事実にはならない。\autocite{src-a07eb39dec04,src-dfe7e0c6ff39}
\end{claimboundary}

\begin{claimboundary}
Claim Boundary: Claude Opus 4.7は保存されたAnthropic newsroomのembedded first-party dataから4月16日のreleaseを確認できる。ただし今回の保存資料だけから詳細なcapabilityやbenchmarkを補完しない。\autocite{src-f6aa679babd7}
\end{claimboundary}

\begin{claimboundary}
Claim Boundary: vendor index／changelogは保存時点で残っている具体的な4月entryだけを扱う。site更新で過去paginationやdetailが失われ得るため、snapshotにない詳細を現在の知識から遡及的に埋めない。\autocite{src-9156c37fb1c2,src-a24fd97eb3d0,src-e547d0d8cbe3,src-eac8dfbeb78c}
\end{claimboundary}

\begin{claimboundary}
Claim Boundary: Transformers release notesはproject chronologyと記載されたsupport変更の一次資料だが、すべてのdeploymentにおけるperformanceやcompatibilityを独立に保証するものではない。\autocite{src-6ddeedd7665f}
\end{claimboundary}
